\documentclass[11pt,a4paper]{article}
\usepackage[margin=0.9in,top=0.8in,bottom=0.85in]{geometry}
\usepackage{fontspec}
\usepackage{xeCJK}
\usepackage{xcolor}
\usepackage{hyperref}
\usepackage{setspace}
\usepackage{titlesec}
\pagestyle{empty}

\setmainfont{Times New Roman}
\setCJKmainfont{Songti SC}
\setCJKsansfont{Heiti SC}
\hypersetup{colorlinks=true,urlcolor=black}

\definecolor{accent}{HTML}{27596A}
\definecolor{softblue}{HTML}{EAF4F8}
\setlength{\parindent}{2em}
\setlength{\parskip}{0.55em}
\linespread{1.18}

\titleformat{\section}{\Large\bfseries\color{accent}\centering}{}{0pt}{}

\begin{document}

\begin{center}
  {\fontsize{22}{28}\selectfont\bfseries 自荐信}\\[8pt]
  {\large 周天昊}\\[4pt]
  天津大学软件工程学院 · 软件工程本科\\
  \href{mailto:zhoutianhao@tju.edu.cn}{zhoutianhao@tju.edu.cn}
\end{center}

\vspace{0.6em}

尊敬的老师/评审专家：

您好！我是天津大学软件工程学院软件工程专业本科生周天昊，计划于 2027 年 6 月毕业。非常感谢您在百忙之中阅读我的自荐材料。我希望申请继续深造或加入与分布式系统、可信计算、区块链共识协议、智能体协作相关的研究团队。过去几年的学习和科研经历让我逐渐明确了自己的方向：以扎实的数学与系统能力为基础，面向真实复杂网络环境中的可靠协作问题，开展可验证、可复现、可落地的研究。

我本科阶段的学习路径具有一定跨学科特点。入学初期，我曾在建筑相关专业学习，并于第 3 学期转入计算机方向。转专业意味着需要在较短时间内补齐大量基础课程，也让我更早意识到自我驱动和学习效率的重要性。转入软件工程后，我系统完成了数学、程序设计、数据结构、算法、操作系统、数据库、计算机系统、软件工程等课程学习，目前 GPA 为 3.82/4.00，平均分 90，专业排名 17/131。核心课程中，自然语言处理 98、计算机系统综合实践 99、设计模式 99、算法设计与分析 96、数据结构 94.8、概率论 95、离散数学 91、形式化方法 92。这些课程训练让我建立了较完整的计算机基础，也为后续阅读论文、实现协议和分析实验结果提供了必要支撑。

科研方面，我目前在天津大学 Tank Lab 参与分布式系统、部分同步/异步协议与 agent 共识相关研究。我的主要工作围绕区块链 BFT 共识、异步 DAG-BFT 协议、部分同步网络模型以及多 agent 协同决策展开。在异步 DAG-BFT 研究中，我参与协议阅读、机制推导、Rust 代码实现与测试、Narwhal-derived 代码库维护、同栈复现实验、WAN/延迟仿真、吞吐与延迟指标分析、图表整理和论文实验撰写支撑。相关工作关注在高传播延迟环境下提升吞吐、确认延迟与 certified block 生成效率，目前一篇 ICDE2027（CCF-A）在投，一篇 IEEE INFOCOM 2027（CCF-A）在投，一篇 ACL 2027（CCF-A）筹备。

在这些科研训练中，我逐渐形成了对系统研究的理解。一个可靠的协议设计需要清晰的问题抽象、严谨的安全性思考、可实现的工程路径和足够稳健的实验验证。阅读协议论文时，我会重点关注网络模型、失效假设、投票与证书机制、leader 选择、提交路径、延迟边界和吞吐瓶颈；实现实验时，我会尽量追踪参数配置、日志处理、指标统计和可复现实验脚本，避免只停留在结论层面。近期我也在思考如何将 BFT 共识机制拓展到多 agent 协作语境，围绕 agent 间状态同步、任务提议、投票确认、异常节点容忍和可验证执行记录建立更可靠的协作流程。这一方向兼具系统安全性、智能体协作和工程实现挑战，也与我长期希望投入的研究主题高度契合。

除主线科研外，我也有跨学科研究经历。转专业前，我曾参与天津大学建筑学院风景园林系王苗团队的社会科学数字化研究，使用 GIS、数据清洗和定量分析方法研究“双减”政策背景下中小学周边开放空间活化潜力，形成科研成果 \textit{Research on the Activate Potential of Open Spaces Around Primary and Secondary Schools Under the Background of the ``Double Reduction'' Policy}。这段经历让我理解到，技术能力的价值常常体现在对真实问题的建模、数据组织和证据表达中。它也让我在后续计算机科研中更重视问题背景、数据质量和结果解释。

工程实践方面，我完成了多个课程与开源实践项目，包括复杂微服务系统与在线评测系统、犯罪率分析与可视化、自然语言处理实验、计算机视觉/计算机可视化课程项目等。其中在线评测系统项目覆盖需求分析、模块实现、接口联调、测试、文档交付与仓库管理；犯罪率分析项目训练了数据预处理、缺失异常处理、统计分析、图表叙事和可复现展示流程；自然语言处理实验覆盖文本清洗、向量化、模型训练、指标评估和误差分析；计算机视觉与可视化项目强化了图像处理、视觉表示和工程展示能力。我熟悉 Python、C++、Go、Rust、TypeScript、Linux、Git 和 LaTeX，也能较熟练地使用 AI-assisted coding 工具辅助代码阅读、实验开发、调试和文档整理。

我珍视团队协作中的长期可靠性。无论在科研项目、课程项目还是学生工作中，我都尽量让自己的工作可追踪、可交接、可复现。曾任学生会综合事务部部长的经历提升了我的沟通与组织能力，也让我在多任务并行时更注重优先级、节奏和责任边界。科研中，我愿意从基础的复现、日志整理、脚本维护做起，也希望在逐步深入问题后提出自己的判断和方案。我相信，真正有价值的研究积累来自持续阅读、反复实现、严格验证和清晰表达。

未来，我希望继续围绕分布式系统、可信计算、区块链共识和 agent 共识开展研究。一方面，我希望进一步提升理论建模与形式化分析能力，理解协议安全性、活性、复杂度与网络假设之间的关系；另一方面，我也希望保持工程实现和实验验证能力，在真实网络延迟、节点规模变化和异常行为条件下评估协议表现。若有机会加入贵团队，我将以认真、稳定和开放的态度投入科研训练，在论文阅读、系统实现、实验复现、结果分析和论文写作中承担扎实工作，并努力成长为能够独立提出问题、实现方案和完成验证的研究者。

再次感谢您阅读我的自荐信。期待有机会进一步交流，也希望能够在未来的学习与研究中贡献自己的努力。

\vspace{1em}
\begin{flushright}
  自荐人：周天昊\\
  2026 年 8 月
\end{flushright}

\end{document}
